On étudie la pile (Cuivre/Aluminium) en utilisant une plaque d'aluminium $A l_{(s)}$ immergée dans une solution aqueuse de chlorure d'aluminium $\left(A l_{(a q)}^{3+}+3 C l_{(a q)}^{-}\right)$, et une plaque de cuivre $C u_{(s)}$ immergée dans une solution aqueuse de sulfate de cuivre II $\left(Cu_{(a q)}^{2+}+SO_{4(a q)}^{2-}\right)$. On relie les deux solutions par un pont salin.

On branche en série avec la pile un conducteur ohmique, un ampèremètre et un interrupteur. On ferme l'interrupteur à un instant $t=0$, l'ampèremètre indique le passage d'un courant électrique d'intensité constante $I=40 ~mA$ pendant la durée de fonctionnement $\Delta t=1 ~h 30 ~min$.

\section*{Données :}
\begin{itemize}
  \item Les deux solutions ont le même volume et la même concentration molaire $C=0,1 ~mol . L^{-1}$;
  \item $F=9,64.10^{4}$ C.mol $^{-1} \quad ; \quad M($ Al $)=27 ~g . mol^{-1}$
  \item La constante d'équilibre associée à l'équation : $3 C u_{(a q)}^{2+}+2 A l_{(s)} \underset{2}{\stackrel{1}{\rightleftarrows}} 3 C u_{(s)}+2 A l_{(a q)}^{3+}$ est $K=10^{200}$.
\end{itemize}

\begin{enumerate}
  \item On se basant sur le critère d'évolution spontanée, déterminer le sens d'évolution du système chimique durant le fonctionnement de la pile.
  \item Déterminer la polarité des électrodes de la pile. Justifier la réponse.
  \item Montrer que l'expression de la masse d'aluminium qui a réagi durant la durée $\Delta t$ est : $m=\frac{I . \Delta t . M(A l)}{3 F}$. Calculer $m$.
\end{enumerate}



